\documentclass[11pt,a4paper]{article}
\usepackage[utf8]{inputenc}
\usepackage[T1]{fontenc}
\usepackage{amsmath}
\usepackage{amsfonts}
\usepackage{amssymb}
\usepackage{graphicx}
\usepackage{booktabs}
\usepackage{xcolor}
\usepackage{hyperref}
\usepackage{geometry}
\usepackage{titlesec}
\usepackage{enumitem}
\usepackage{fancyhdr}

% Geometry setup
\geometry{margin=1in}

% Colors
\definecolor{agrigreen}{HTML}{059669}
\definecolor{agriamber}{HTML}{D97706}
\definecolor{agriblue}{HTML}{2563EB}

% Title Formatting
\titleformat{\section}{\large\bfseries\color{agriblue}}{}{0em}{}[\titlerule]
\titleformat{\subsection}{\bfseries\color{agrigreen}}{}{0em}{}

% Header/Footer
\pagestyle{fancy}
\fancyhf{}
\rhead{\textit{AgriSense Intelligence Hub}}
\lhead{Executive Summary}
\cfoot{\thepage}

\begin{document}

\begin{center}
    {\LARGE \textbf{AgriSense Intelligence Hub: Executive Summary}} \\
    \vspace{0.2cm}
    {\large \textit{Hybrid AI Framework for Precision Agriculture \& Market Intelligence}} \\
    \vspace{0.5cm}
    \textbf{Status: Production-Ready (Phase 3 Finalized)}
\end{center}

\section{Project Overview}
AgriSense AI is a unified intelligence layer designed to mitigate financial risk and biological uncertainty in Indian agriculture. By utilizing a hybrid multi-model architecture, the platform bridges the gap between commodity price forecasting, soil-crop suitability, and early-warning disease detection.

\section{Core Performance Metrics}
The following table summarizes the validated performance of the AgriSense production modules as of May 5th, 2026.

\begin{table}[h!]
\centering
\begin{tabular}{@{}llll@{}}
\toprule
\textbf{Module} & \textbf{Core Model} & \textbf{Primary Metric} & \textbf{Value} \\ \midrule
\textbf{Price (Primary)} & Temporal Fusion Transformer (TFT) & Accuracy (1-SMAPE) & \textbf{96.15\%} \\
\textbf{Price (Secondary)} & LightGBM (GOSS/EFB Optimized) & SMAPE & \textbf{8.80\%} \\
\textbf{Crop Intel} & Leaf-wise GBDT Classifier & F1-Score & \textbf{0.99} \\
\textbf{Disease Radar} & EfficientNet-B0 (Transfer Learning) & Accuracy & \textbf{98.42\%} \\ \bottomrule
\end{tabular}
\caption{Production Model Benchmarks}
\end{table}

\section{The Dual-Track Forecasting Strategy}
To capture both long-term market trends and short-term volatility, the system employs a synchronized forecasting track:
\begin{itemize}[leftmargin=*]
    \item \textbf{TFT Track:} Captures long-range dependencies across 30-day context windows, providing quantile-based risk intervals for financial planning.
    \item \textbf{LightGBM Track:} Provides ultra-low latency point predictions for 14-day market timing, leveraging over 70 engineered temporal features including Fourier seasonality and NASA weather shocks.
\end{itemize}

\section{Architectural Integrity}
The system is built on a \textbf{Decoupled Engine Pattern}. The \texttt{Engines/} directory provides a canonical interface for all inference logic, ensuring the Next.js frontend can interact with high-compute PyTorch models and fast XGBoost/LGBM classifiers through a unified Flask-based Intelligence Hub.

\section{Technocratic Assessment}
While the \textbf{Crop Recommendation} and \textbf{Disease Detection} modules exhibit near-perfect benchmark performance, the system frames these as \textit{Biological Feasibility Indicators}. This honest assessment acknowledges the "Data-Sparse" constraints of specialized datasets (2,200 records) while emphasizing the high fidelity of the "Data-Rich" Mandi Price track (250,000+ records).

\section{Conclusion}
AgriSense AI represents a significant leap in data-driven agriculture, providing farmers with a robust toolkit for profit maximization and risk mitigation. The project is fully documented, modular, and ready for collaborative scaling.

\vfill
\begin{center}
    \textit{Developed by Karthik Reddy (230035) \& Akash G (230098) --- Rishihood University}
\end{center}

\end{document}
